\section{\texorpdfstring{\(3\times3\) 행렬곱의 rank (\(\mathbb F_2\))}{3×3 행렬곱의 rank (F2)}}

2026-10-12. 결과: 22를 찾지도, 배제하지도 못함. 다음 공격 지점은 정해
둠.

\subsection{문제}

\(3\times3\) 행렬 두 개를 곱할 때 필요한 곱셈의 최소 횟수(덧셈은 공짜).
쌍선형 알고리즘

\[
m_t=\Big(\sum_{i,j}u^{(t)}_{ij}a_{ij}\Big)\Big(\sum_{i,j}v^{(t)}_{ij}b_{ij}\Big),\qquad c_{ij}=\sum_tw^{(t)}_{ij}m_t
\]

이 가능한 최소 항 수가 행렬곱 텐서 \(\langle3,3,3\rangle\)의 rank
\(R\)이다.

\begin{itemize}
\tightlist
\item
  교과서 방식 27, Laderman(1976) 23이 지금도 최선. SAT 탐색으로 23항
  해가 수천 개 나왔지만 22는 없음.
\item
  \(2\times2\)는 Strassen 7, 그리고 7이 최소(Winograd).
\end{itemize}

\subsubsection{왜 21이 분수령인가}

\[
\log_27\approx2.807,\quad\log_323\approx2.854,\quad\log_322\approx2.814,\quad\log_321\approx2.771
\]

\(R\le21\)이어야 \(3\times3\) 재귀가 Strassen을 이긴다.

\subsubsection{현재 하한}

\begin{itemize}
\tightlist
\item
  일반 체: Bläser(2003) 19
\item
  \(\mathbb F_2\): Chengu Wang(2026)이 \(19\to20\)
  (\href{https://arxiv.org/abs/2603.07280}{arXiv 2603.07280}), 다시
  \(20\to21\) (\href{https://arxiv.org/abs/2609.06725}{arXiv
  2609.06725}). 따라서 \(\mathbb F_2\)에서 \(21\le R\le23\).
\end{itemize}

공학적 의미는 크지 않다. 실제 GPU·시스톨릭 배열은 교과서식 \(n^3\)을
쓴다. 영향은 곱셈이 비싼 환경(큰 정수, 암호, FPGA DSP 수 제약)과 행렬곱
지수 \(\omega\) 연구 정도.

\subsection{시도 1: flip graph 탐색 (상한 쪽)}

Kauers--Moosbauer(2022) 방식. 27항 분해에서 출발해, 두 항이 인자 하나를
공유하면 다른 인자를 섞는 flip을 반복하고, 항이 0이 되거나 합쳐지면
rank가 떨어진다. 찾은 해는 매번 텐서를 다시 계산해 검증.

\begin{itemize}
\tightlist
\item
  코드: \url{code/flip333.c}
  (\texttt{gcc\ -O3\ -o\ flip333\ flip333.c\ \&\&\ ./flip333\ \textless{}seed\textgreater{}\ \textless{}초\textgreater{}\ \textless{}한도\textgreater{}})
\item
  규모: 코어 2개 × 9분, 약 \(10.4\times10^8\) flip
\item
  결과: 독립적으로 도달한 rank 23 해 87개, 각각 23에서 약 1200만 걸음 더
  걸었지만 22로 떨어진 경우 0번
\end{itemize}

원 논문보다 훨씬 작은 규모라 ``22가 없다''는 증거로는 약하다.

\subsection{시도 2: Wang의 하한 증명 확장 (하한 쪽)}

Wang의 증명 구조:

\begin{enumerate}
\def\labelenumi{\arabic{enumi}.}
\tightlist
\item
  \textbf{부분공간 하한표.} 첫 입력을 부분공간 \(S\)로 제한한 텐서의
  rank 하한 \(L(S)\)를 전부 계산.
\item
  \textbf{용량 논증.} \(r\)항 분해가 있다면 첫 인자가 \(E=S^\perp\) 안에
  있는 항은 \(r-L(S)\)개 이하. 이를 만족하는 첫 인자 배치(프로파일)를
  전수 탐색해 \(r=20\)에서 0개 → \(R\ge21\).
\end{enumerate}

22를 얻으려면 \(r=21\) 프로파일이 없음을 보여야 한다. \(r\)이 1 오르면
모든 용량이 1씩 늘고, \(r=20\)에서 쓴 핵심 보조정리(두 행렬의 합이 rank
1)가 \(r=21\)에서는 성립하지 않아 그대로는 전수 탐색이 안 된다.

한 것:

\begin{itemize}
\tightlist
\item
  용량 판정기를 C로 독립 재구현: \url{code/capacity.c}, 인증서 변환
  \url{code/parse_cert.py}. 인증서는
  \href{https://github.com/wcgbg/matrix-multiplication-n333r21f2}{wcgbg/matrix-multiplication-n333r21f2}의
  \texttt{certs/matrix/cert\_matrix\_q02\_n333.pb.txt} (Git LFS로 받아야
  함). 부분공간 8,283,458개, 1차원 \(E\)에서 \(L=19\)로 논문과 일치.
  실제 23항 해 20개가 모두 \(r=23\) 조건을 통과(양성 대조).
\item
  \textbf{실험 1:} 23항 해 20개에서 항 2개를 뺀 21항 배치 5,060개 →
  \(r=21\) 통과 0개. 1개를 뺀 22항 배치 460개 → \(r=22\) 통과 0개.
  해마다 용량이 꽉 찬 부분공간이 33\textasciitilde142개라서 알려진 해
  근처에는 \(r=21\) 프로파일이 없다.
\item
  \textbf{실험 2:} 처음부터 담금질 탐색. \(r=21\)에서 위반량 5까지,
  대조군 \(r=23\)도 1에서 멈춤 → 탐색기가 약하다는 것만 보여줌. 결론
  없음.
\end{itemize}

\subsubsection{판정기 실행}

\begin{verbatim}
cp <wcgbg 저장소>/certs/matrix/cert_matrix_q02_n333.pb.txt .
python3 parse_cert.py                 # -> orbits.txt (궤도 대표와 하한 L)
gcc -O3 -o capacity capacity.c -lm
./capacity control r23.txt            # 23항 해의 첫 인자가 r=23 용량 조건을 통과하는지 (양성 대조)
./capacity sub r23.txt                # 항 2개/1개를 뺀 배치를 r=21/22로 검사 (실험 1)
./capacity sa 21 400000 1             # r=21 담금질 탐색 (실험 2)
\end{verbatim}

\texttt{r23.txt}는 \texttt{flip333}이 rank 23에 도달하면 저장하는
파일이다. 판정기는 부분공간 표를 해시로 펼치는 데 메모리 약 600MB를
쓴다.

\subsection{다음 공격 지점}

1차원 \(E=\langle U\rangle\), 즉 hyperplane 제한의 bound가 세
궤도(\(U\)의 rank 1, 2, 3) 모두

\[
19\le R_{\mathbb F_2}\big(T_{\ker U}\big)\le21
\]

이다. 이걸 20으로 올리면 \(r=21\)에서 1차원 용량이 1이 되어 ``프로파일
원소는 서로 다르다''는 성질이 살아나고, \(r=21\) 탐색을 잘라낼 첫 단추가
된다. 부분문제는 \(8\times9\times9\) 텐서이고, Wang이 20→21을 한 구조를
한 단계 작은 크기에서 반복하는 셈이다. (2차원 \(E\) 일부는 \(L=U=18\)로
이미 확정이라 더 못 올린다.)

순서: \(T_{\ker U}\) 세 개에 \(r=19\) 프로파일 탐색 → \(R\ge20\) 시도 →
성공하면 \(r=21\)용 새 구조 보조정리 찾기.
